\subsection{Aproximação curta e região de captura}
\label{sec:aprox_curta_regiao_captura}

A fase de aproximação curta tem início quando o perseguidor atinge a vizinhança imediata do alvo.
A partir desse momento, a estratégia de aproximação passa a ser regida por requisitos mais restritos de posição, velocidade relativa e atitude, definidos em torno da porta de acoplamento do alvo.
A região ao redor da porta de acoplamento é delimitada por valores máximos admissíveis de posição e velocidade relativa em cada eixo do \gls{lvlh}.

Do ponto de vista translacional, a aproximação curta é considerada bem-sucedida quando a posição relativa do perseguidor encontra-se no interior dos limites impostos de posição e velocidade para cada eixo do \gls{lvlh}.
Esse conjunto de condições permite que o manipulador robótico possa ser acionado para realizar o alinhamento fino e a captura mecânica, sem exigir grandes correções de translação por parte da base.
